\corrigesChapitre{10}{Échantillonnage et décision statistique}
\begin{corrige}{1}{Fréquence et proportion}
\begin{questions}
\item $E(X)=40$ et $E(F)=0{,}4$.
\item $\sqrt{0{,}4\times0{,}6/100}\simeq0{,}04899$.
\item $f=0{,}45$, différente de $0{,}4$ ; cette différence peut résulter de la fluctuation.
\end{questions}
\end{corrige}
\begin{corrige}{2}{Effet de la taille}
\begin{questions}
\item Par $\sqrt4=2$.
\item Par $\sqrt{100}=10$.
\item Non : la formule traite la fluctuation dans un modèle d’échantillonnage, pas une erreur systématique de recrutement.
\end{questions}
\end{corrige}
\begin{corrige}{3}{Intervalle exact simple}
\begin{questions}
\item $1-2(1/16)=7/8=0{,}875$.
\item Non : $0{,}875<0{,}95$.
\item Il faut ici inclure les deux extrêmes pour dépasser $0{,}95$ ; $[0,4]$ a probabilité $1$.
\end{questions}
\end{corrige}
\begin{corrige}{4}{Complément}
\begin{questions}
\item $1-0{,}96=0{,}04$.
\item Oui, $0{,}04<0{,}05$.
\item Non : rare ne signifie pas impossible.
\end{questions}
\end{corrige}
\begin{corrige}{5}{Tester une pièce}
\begin{questions}
\item $P(X\leq1)+P(X\geq9)=2(1+10)/1024=22/1024\simeq0{,}021484$.
\item $9\notin I$ : on rejette l’hypothèse selon cette règle.
\item $6\in I$ : on ne rejette pas. Cela ne prouve pas que la pièce est équilibrée.
\end{questions}
\end{corrige}
\begin{corrige}{6}{Une borne de capacité}
\begin{questions}
\item $0{,}2^4+4(0{,}8)(0{,}2)^3+6(0{,}8)^2(0{,}2)^2=0{,}1808$.
\item $1-0{,}8^4=0{,}5904$.
\item Il faut un cumul au moins $0{,}5$ : la plus petite capacité est $3$, puisque le cumul en $2$ est inférieur.
\end{questions}
\end{corrige}
\begin{corrige}{7}{Simulation reproductible}
\begin{questions}
\item Elle initialise un générateur pseudo-aléatoire ; la même graine reproduit la même suite de tirages avec la même implémentation.
\item On convertit le nombre de succès en fréquence observée.
\item Une fréquence par échantillon ; sa dispersion illustre la fluctuation d’échantillonnage.
\end{questions}
\end{corrige}
\begin{corrige}{8}{Sondage biaisé}
\begin{questions}
\item Non : la taille ne suffit pas à assurer la qualité du recrutement.
\item Les visiteurs ne représentent pas nécessairement la population visée ; les personnes motivées répondent davantage. Le libellé peut aussi influencer les réponses.
\item Un plan de sondage aléatoire maîtrisé permet de modéliser le recrutement ; il faut encore traiter les non-réponses et la qualité de mesure.
\end{questions}
\end{corrige}
\begin{corrige}{9}{Deux types de risque}
\begin{questions}
\item Le risque de rejeter l’hypothèse alors qu’elle est vraie, dans le modèle annoncé.
\item Non : la probabilité de l’hypothèse après observation n’est pas fournie par cette procédure.
\item Une autre valeur du paramètre peut produire une observation compatible avec la zone ; le test peut manquer de puissance pour la distinguer.
\end{questions}
\end{corrige}
\begin{corrige}{10}{Moyennes d’un dé}
\begin{questions}
\item $\sqrt{35/12}/10\simeq0{,}170783$.
\item $0{,}341565$ environ.
\item Pour chaque moyenne $m$, tester $|m-3{,}5|\leq0{,}341565$ ; compter les succès puis diviser par le nombre de moyennes simulées. On conserve davantage de chiffres dans le calcul effectif.
\end{questions}
\end{corrige}
\begin{corrige}{11}{Une zone rare}
\begin{questions}
\item $(1+10+10+1)/1024=11/512\simeq0{,}021484$.
\item Oui : $0{,}021484<0{,}025$.
\item C’est son complément sur l’ensemble des valeurs entières possibles. La probabilité de $[2,8]$ dépasse donc $0{,}975$.
\end{questions}
\end{corrige}
\begin{corrige}{12}{Construire plutôt que deviner}
\begin{questions}
\item Pour $r=0,1,\ldots,10$, calculer $\sum_{k=10-r}^{10+r}\binom{20}k2^{-20}$ et s’arrêter au premier cumul au moins $0{,}95$.
\item $P(6\leq X\leq14)=0{,}9586105347$ ; $P(5\leq X\leq15)=0{,}9881820679$.
\item Le rayon $4$ convient. Le rayon $3$ donne $P(7\leq X\leq13)\simeq0{,}884682<0{,}95$ : $4$ est minimal.
\end{questions}
\end{corrige}
